\chapter{La première flèche}
\label{chap:premiere-fleche}

L'intérêt de Kael pour l'arc avait commencé bien avant sa rencontre avec Robin
Mangil. Kamatsu et lui avaient un jour fabriqué ce qu'ils considéraient très
généreusement comme un arc. Ils avaient choisi une branche suffisamment souple
pour accepter d'être courbée, récupéré un morceau de corde et taillé quelques
projectiles qui, avec beaucoup de bonne volonté, pouvaient passer pour des
flèches.

L'ensemble fonctionnait. Pas très bien, certes, mais suffisamment pour
permettre à une flèche de parcourir quelques mètres avant de tomber dans
l'herbe, ce qui leur avait semblé être une réussite incontestable. Ils avaient
recommencé encore et encore, et Kael avait rapidement pris goût à l'exercice.
Il aimait cette étrange sensation qui précédait le tir : choisir un point,
tendre la corde, retenir son geste pendant un instant, puis libérer la flèche
et découvrir si ce qu'il avait imaginé correspondait à ce qui allait réellement
se produire.

\scenebreak

Robin Mangil voyageait avec la caravane depuis quelques jours lorsque Kael
commença à remarquer ses absences. Lors des haltes, le demi-elfe s'éloignait
régulièrement des chariots avec son arc et son carquois. Il revenait quelque
temps plus tard, souvent sans que personne semble particulièrement s'être
préoccupé de ce qu'il était parti faire.

Kael, lui, s'en préoccupait.

Un soir, il le suivit à distance et le vit s'arrêter à l'écart du camp pour
choisir un arbre. Robin tira une première flèche, puis recommença. Kael resta
suffisamment loin pour ne pas le déranger et observa. Les gestes semblaient
simples : Robin prenait une flèche, la plaçait sur la corde, levait son arc et
tirait.

Pourtant, plus Kael regardait, plus il remarquait de petites différences dans
la position de ses pieds, la manière dont son corps accompagnait le mouvement,
le temps qu'il prenait avant certains tirs ou encore la hauteur de l'arc
lorsque la cible était plus éloignée. Kael revint le lendemain, puis le jour
suivant.

La troisième fois, Robin abaissa son arc sans tirer.

— Tu comptes rester caché encore longtemps ?

Kael sortit de derrière l'arbre depuis lequel il observait, sans sembler
particulièrement embarrassé d'avoir été découvert. Robin le regarda approcher.

— Qu'est-ce que tu regardes ?

— Vos mains.

Robin baissa les yeux vers elles.

— Mes mains ?

— Vous ne tenez pas l'arc pareil quand vous tirez loin.

Robin observa le garçon avec un peu plus d'attention.

— Tu as déjà tiré ?

— Oui.

— Avec quoi ?

Kael hésita.

— Un arc.

— J'avais deviné.

Kael réfléchit quelques secondes.

— Je peux vous montrer.

\scenebreak

Le lendemain, Kael revint avec l'arc qu'il avait fabriqué avec Kamatsu et le
tendit fièrement à Robin. Le demi-elfe examina la branche, la corde, puis les
quelques flèches improvisées que Kael avait apportées avec lui. Il retourna
l'arc entre ses mains avant de regarder le garçon.

— Et vous avez vraiment tiré avec ça ?

— Oui.

— Plusieurs fois ?

— Oui.

Robin examina de nouveau l'objet.

— Et personne n'a perdu un œil ?

— Non.

Robin lui rendit l'arc.

— Alors c'est déjà une réussite.

Kael ne sut pas exactement s'il s'agissait d'un compliment. Robin, lui, avait
déjà pris son propre arc.

— Tiens.

Kael le saisit et sentit immédiatement la différence. L'arme était plus lourde
que la branche qu'il utilisait habituellement, mais surtout la corde lui opposa
une résistance qu'il n'avait pas prévue. Lorsqu'il essaya de la tirer, elle
bougea à peine.

Robin sourit.

— Toujours envie d'essayer ?

Kael raffermit sa prise et tira de nouveau. Cette fois, la corde recula,
difficilement. Robin lui tendit une flèche.

— Ne tire pas tout de suite.

Kael plaça la flèche comme il l'avait fait avec les siennes.

— Tes pieds.

Kael baissa les yeux.

— Quoi, mes pieds ?

Robin les repositionna, puis corrigea légèrement ses épaules.

— Ton bras.

Kael obéit.

— Et maintenant ?

— Maintenant, tu regardes ta cible.

Kael fixa l'arbre devant lui, tendit la corde aussi loin qu'il le pouvait et
tira. La flèche passa largement à côté du point qu'il visait, mais se planta
malgré tout dans le tronc. Kael regarda successivement la flèche et l'endroit
qu'il avait voulu atteindre.

— J'ai raté.

— Oui. Mais c'est ta première flèche avec un véritable arc, celui-ci est
beaucoup trop grand pour toi et tu arrives à peine à le tendre.

Robin récupéra une autre flèche et la lui tendit.

— Alors, pour une première : pas mal.

Kael sourit.

— Je peux recommencer ?

Robin désigna ses jambes.

— D'abord, tes pieds.

Ce fut le début de leurs entraînements.

\scenebreak

Robin découvrit rapidement que Kael avait des facilités. Il manquait de
technique et d'expérience, mais observait beaucoup et reproduisait rapidement
ce qu'on lui montrait. Certaines corrections ne devaient lui être répétées que
quelques fois avant de devenir naturelles.

D'autres beaucoup plus souvent.

— Ton coude.

Kael le corrigea et tira.

— Ton coude.

Kael reprit une flèche.

— Encore ?

— Encore.

Kael soupira.

Robin lui apprit à placer ses pieds, à positionner son corps et à ne pas lutter
inutilement contre son arc. Il corrigea sa manière de tenir la corde, puis
celle de viser, mais surtout il lui apprit à regarder. Pas seulement sa cible :
la distance, le vent, le terrain, les branches qui pouvaient gêner le passage
d'une flèche et parfois même ce que la cible allait faire avant qu'elle ne le
fasse.

Lorsqu'un passage en ville permit aux voyageurs de refaire leurs provisions,
Kael demanda à Robin de l'accompagner chez un marchand. Il avait économisé
presque tout l'argent qu'il possédait depuis qu'il vivait avec la caravane et
voulait acheter son premier véritable arc.

\scenebreak

Kael choisit naturellement l'un des plus grands.

— Non, dit immédiatement Robin.

Kael tourna la tête.

— Pourquoi ?

— Parce qu'il est trop grand.

— Je peux le tendre.

— Probablement.

Kael attendit la suite. Robin reprit l'arc et le remit à sa place.

— Tu veux un arc que tu peux utiliser, pas un arc qui donne l'impression que
tu sais l'utiliser.

Kael regarda l'arme, puis Robin, puis de nouveau l'arme.

— Non.

Ils continuèrent donc à chercher. Plusieurs arcs passèrent entre les mains de
Kael avant que Robin n'en sorte un autre. Il était d'occasion : le bois portait
quelques marques, la poignée avait déjà beaucoup servi et il n'avait
certainement rien d'impressionnant à côté de celui du demi-elfe.

Mais il était solide. Et surtout, lorsque Kael tendit la corde, il put le faire
sans que tout son corps participe à l'effort.

Robin acquiesça.

— Celui-là.

Kael regarda le prix, puis sa bourse, et compta ses pièces deux fois.
Lorsqu'il ressortit de la boutique, il ne lui restait presque plus rien.

Mais il avait son arc.

\scenebreak

À partir de ce jour, une habitude s'installa. Lors des haltes suffisamment
longues, Robin prenait son arc, Kael prenait le sien et tous deux partaient
s'entraîner. Les distances augmentèrent et les cibles diminuèrent. Kael apprit
à tirer debout, accroupi, après avoir couru ou lorsqu'il ne disposait que de
quelques secondes pour préparer sa flèche.

Un soir, Robin lui désigna une branche qui oscillait avec le vent.

— Celle-là.

Kael observa la branche.

— Elle bouge.

— C'est le principe.

— C'est trop difficile.

Robin tourna légèrement la tête vers lui.

— Ah bon ?

Kael le regarda.

— Tu n'en es pas capable ?

Kael ne répondit rien. Il prit simplement une flèche.

Robin eut un léger sourire, mais se garda bien d'ajouter quoi que ce soit.
Kael leva son arc et tira. La flèche passa derrière la branche. Il en prit
aussitôt une autre ; cette fois, elle passa devant.

Kael encocha une troisième flèche, attendit un peu plus longtemps et tira de
nouveau. Nouvel échec. Il cherchait déjà une quatrième flèche lorsque Robin
intervint.

— Arrête de tirer là où elle est.

Kael se tourna vers lui. Robin désigna la branche.

— Tire là où elle sera.

Kael fronça les sourcils et regarda de nouveau sa cible. La branche oscillait
avec le vent, partait dans une direction, ralentissait, puis revenait avant de
recommencer. Cette fois, il ne chercha pas à la suivre avec son arc. Il observa
son mouvement pendant quelques instants, essaya d'en comprendre le rythme, puis
encocha une nouvelle flèche.

Il attendit.

La branche repartit. Kael leva son arc et tira. La flèche traversa l'endroit où
la branche n'était pas encore. Une fraction de seconde plus tard, celle-ci
revint sur sa trajectoire.

La flèche la frappa.

La branche se brisa.

Kael resta quelques secondes à regarder le résultat. Robin sourit.

— Maintenant, tu commences à tirer.

\scenebreak

Lorsque vint le moment pour Robin de reprendre sa propre route, il ne donna à
Kael ni arme ni souvenir particulier. Kael avait déjà son arc, et Robin lui
avait donné quelque chose de plus utile. Avant de partir, le demi-elfe désigna
simplement les yeux de Kael, puis les alentours.

— Et avant de tirer ?

Kael connaissait désormais parfaitement la réponse.

— Je regarde.

Robin acquiesça.

— Alors ça devrait aller.

Le lendemain, Kael sortit son arc comme d'habitude. Il n'y en avait plus deux,
mais il alla s'entraîner quand même.

\scenebreak

Avec le temps, Kael commença à utiliser son arc autrement que pour atteindre
des morceaux de bois.

Il essaya de chasser.

Ses premières tentatives furent surtout de longues promenades dont il revenait
les mains vides. Bartiméus prenait alors parfois le temps de venir l'accueillir,
de renifler son carquois vide puis de repartir avec suffisamment de dignité
pour convaincre Kael que le chat se moquait de lui.

Kael repartait le lendemain.

Il apprit rapidement que toucher une cible immobile et atteindre un animal
étaient deux choses très différentes. Il fallait trouver sa trace, l'approcher
sans faire de bruit, tenir compte du vent et parfois attendre longtemps avant
de pouvoir seulement envisager de tirer. Sur ce point, certaines choses
apprises dans son enfance lui revinrent naturellement.

Puis, un soir, Kael revint enfin au camp avec un petit animal. Bartiméus fut le
premier à venir inspecter sa prise.

— Et non, ce n'est pas pour toi.

Le chat miaula.

— C'est pour tout le monde.

Nouveau miaulement.

— Tu en auras un morceau.

Bartiméus sembla considérer cette proposition comme plus raisonnable et
repartit vers le camp.

L'animal fut préparé et partagé autour du feu. Chacun n'en reçut qu'un peu,
Bartiméus compris. Kael mangea sa part avec les autres, mais regarda surtout
les membres de la caravane partager ce qu'il avait rapporté. Il ne s'agissait
que d'un petit animal, à peine de quoi améliorer le repas de quelques
personnes, et pourtant il souriait.

Pour la première fois, ce qu'il avait appris sur les routes lui avait permis
de rapporter quelque chose à la famille qui l'avait accueilli.